% ============================================================
% PHẦN 1. GIỚI THIỆU
% ============================================================
\setroadmap{-1}
\section{Giới thiệu}

% ------------------------------------------------------------
% F1.1 Động lực
% ------------------------------------------------------------
\begin{frame}
\frametitle{LLM trả lời trôi chảy nhưng chưa đủ căn cứ pháp lý}
\begin{columns}[T,onlytextwidth]
\column{0.55\textwidth}
\small
\begin{itemize}
  \item LLM làm tốt nhiều bài toán hỏi đáp, nhưng hỏi đáp pháp luật
        đáng tin cậy cần \textbf{đồng thời}:
  \begin{itemize}
    \item căn cứ luật thành văn phù hợp;
    \item chuỗi suy luận truy vết được: điều kiện $\rightarrow$
          sự kiện trung gian $\rightarrow$ hệ quả pháp lý.
  \end{itemize}
  \item RAG bổ sung bằng chứng bên ngoài trước bước sinh
        {\color{schemaGray}(Lewis và cs., 2020; Gao và cs., 2024)}.
  \item Nhưng truy xuất ngữ nghĩa thường xem các đoạn văn là
        \alert{độc lập}, không biểu diễn tường minh quan hệ
        điều kiện--hệ quả nhiều bước.
\end{itemize}

\column{0.41\textwidth}
\centering
\begin{tikzpicture}[font=\scriptsize]
  % RAG cơ bản
  \node[font=\scriptsize\bfseries,text=schemaGray,anchor=west] at (-0.2,1.7) {RAG cơ bản};
  \node[minibox=schemaGray,minimum width=0.95cm] (q) at (0.35,0.65) {Câu\\hỏi};
  \node[minibox=schemaGray,minimum height=0.34cm,minimum width=1.05cm] (d1) at (1.85,1.12) {đoạn 1};
  \node[minibox=schemaGray,minimum height=0.34cm,minimum width=1.05cm] (d2) at (1.85,0.65) {đoạn 2};
  \node[minibox=schemaGray,minimum height=0.34cm,minimum width=1.05cm] (d3) at (1.85,0.18) {đoạn 3};
  \node[fit=(d1)(d3),inner sep=1.5pt] (docs) {};
  \node[minibox=schemaGray,minimum width=0.8cm] (llm) at (3.3,0.65) {LLM};
  \node[minibox=schemaGray,minimum width=1.05cm] (ans) at (4.7,0.65) {Câu\\trả lời};
  \draw[flow] (q) -- (docs);
  \draw[flow] (docs) -- (llm);
  \draw[flow] (llm) -- (ans);
  \node[tag] at (1.85,-0.3) {các đoạn độc lập};
  % Điều bài toán cần
  \node[font=\scriptsize\bfseries,text=stepGraph!80!black,anchor=west] at (-0.2,-0.85) {Điều bài toán cần};
  \node[minibox=stepGraph] (c) at (0.35,-1.75) {Điều\\kiện};
  \node[minibox=stepGraph,double,double distance=0.8pt] (m) at (2.5,-1.75) {Sự kiện\\trung gian};
  \node[minibox=stepGraph] (e) at (4.6,-1.75) {Hệ quả\\pháp lý};
  \draw[causes=stepGraph] (c) -- node[above=1pt,font=\scriptsize]{Điều $a$} (m);
  \draw[causes=stepGraph] (m) -- node[above=1pt,font=\scriptsize]{Điều $b$} (e);
  \node[tag] at (2.5,-2.5) {quan hệ có hướng, mang nguồn điều luật};
\end{tikzpicture}
\end{columns}

\vspace{0.35cm}
\keymsg{Truy xuất đúng tài liệu chưa chắc dẫn đến suy luận đúng.}
\end{frame}

% ------------------------------------------------------------
% F1.2 Ví dụ xuyên suốt
% ------------------------------------------------------------
\begin{frame}
\frametitle{Ví dụ xuyên suốt: câu hỏi cần nối hai điều luật}
\vspace{-0.2cm}
\begin{block}{Câu hỏi \exid{} (suy luận ngược, 2 bước)}
\small ``Kết quả \textit{Có tình tiết giảm nhẹ trách nhiệm hình sự} trong chain này bắt
nguồn từ sự kiện nào trước khi đi qua bước \textit{Vượt quá giới hạn phòng vệ chính đáng}?''
\end{block}

\vspace{0.05cm}
\centering
\begin{tikzpicture}
  \node[evnode=stepGraph,text width=2.75cm,dashed,line width=0.9pt] (a) at (0,0) {Chống trả rõ ràng\\quá mức cần thiết};
  \node[bridgenode=stepGraph,text width=2.75cm] (b) at (5.55,0) {Vượt quá giới hạn\\phòng vệ chính đáng};
  \node[evnode=stepGraph,text width=2.75cm] (c) at (11.1,0) {Có tình tiết giảm nhẹ\\trách nhiệm hình sự};
  \draw[causes=stepGraph] (a) -- node[above=1pt,font=\scriptsize]{rule 47 $\cdot$ Điều 22} (b);
  \draw[causes=stepGraph] (b) -- node[above=1pt,font=\scriptsize]{rule 106 $\cdot$ Điều 51} (c);
  \node[tag,below=1pt of a] {\textbf{đáp án cần tìm}};
  \node[tag,below=1pt of b] {sự kiện cầu nối (phải đi qua)};
  \node[tag,below=1pt of c] {điểm xuất phát của câu hỏi};
\end{tikzpicture}

\vspace{-0.05cm}
\raggedright
\begin{columns}[T,onlytextwidth]
\column{0.55\textwidth}
\begin{block}{Vì sao khó?}
\footnotesize
\begin{itemize}
  \item Hỏi ngược: từ hệ quả về điều kiện ban đầu.
  \item Phải nối hai điều luật: Điều 22 và Điều 51.
  \item Sự kiện cầu nối không là đáp án nhưng phải đi qua.
\end{itemize}
\end{block}
\column{0.41\textwidth}
\begin{alertblock}{Closed-book (abstain) trả lời}
\footnotesize ``Không đủ thông tin để xác định sự kiện trước bước
`Vượt quá giới hạn PVCĐ'.'' $\Rightarrow$ \labUnc
\end{alertblock}
\end{columns}

\vspace{0.08cm}
{\scriptsize\color{schemaGray} Viết tắt: PVCĐ -- phòng vệ chính đáng; TNHS -- trách nhiệm hình sự; BLHS -- Bộ luật Hình sự.}
\end{frame}

% ------------------------------------------------------------
% F1.3 Phát biểu bài toán
% ------------------------------------------------------------
\begin{frame}
\frametitle{Bài toán: quyết định, câu trả lời và căn cứ truy vết được}
\vspace{-0.2cm}
\begin{columns}[T,onlytextwidth]
\column{0.48\textwidth}
\begin{block}{Đầu vào}
\footnotesize
\begin{itemize}
  \item Câu hỏi pháp lý $q$.
  \item Đồ thị có hướng $G=(V,E)$: nút là sự kiện pháp lý, cạnh là rule.
  \item Tập rule $\mathcal R$ xác định ngữ nghĩa suy luận.
\end{itemize}
\end{block}
\column{0.48\textwidth}
\begin{block}{Đầu ra}
\footnotesize
\begin{itemize}
  \item Đường suy luận $P^{*}=v_0\xrightarrow{r_1}v_1\cdots\xrightarrow{r_h}v_h$.
  \item Quyết định $z$ và tập bằng chứng $\mathcal E^{*}$.
  \item Câu trả lời $a$ và căn cứ trích dẫn $C$.
\end{itemize}
\end{block}
\end{columns}

\vspace{0.05cm}
\begin{center}
$f:\ (q,G,\mathcal R)\ \longrightarrow\ (P^{*},\,z,\,\mathcal E^{*},\,a,\,C)$\\[1pt]
{\footnotesize\dien hệ thống phải trả về chuỗi lập luận và quyết định đã kiểm chứng,
không chỉ một câu trả lời.}
\end{center}

\vspace{-0.2cm}
{\footnotesize
\renewcommand{\arraystretch}{1.05}
\begin{tabular}{@{}l p{0.72\textwidth}@{}}
\toprule
\labSup & Luật (qua đồ thị và rule) hỗ trợ quan hệ được hỏi.\\
\labRej & Bác bỏ khẳng định ``trực tiếp'' hoặc ``bỏ qua được bước trung gian''
          khi luật chỉ hỗ trợ chuỗi gián tiếp $A\to B\to C$.\\
\labUnc & Không đủ căn cứ để kết luận (từ chối trả lời).\\
\bottomrule
\end{tabular}}

\vspace{0.05cm}
{\scriptsize\color{schemaRed} Lưu ý: cạnh ``nhân quả'' là phụ thuộc quy phạm trích từ luật, không ước lượng từ dữ liệu quan sát.}
\end{frame}

% ------------------------------------------------------------
% F1.4 Nghiên cứu liên quan
% ------------------------------------------------------------
\begin{frame}
\frametitle{Nghiên cứu liên quan: RAG, đồ thị và suy luận quy tắc}
\vspace{-0.05cm}
{\scriptsize
\setlength{\tabcolsep}{3pt}
\renewcommand{\arraystretch}{1.0}
\begin{tabular}{@{}>{\raggedright\arraybackslash}p{0.39\textwidth}
                >{\raggedright\arraybackslash}p{0.28\textwidth}
                >{\raggedright\arraybackslash}p{0.29\textwidth}@{}}
\toprule
\textbf{Hướng nghiên cứu} & \textbf{Ý tưởng chính} & \textbf{Hạn chế với QA pháp luật đa bước}\\
\midrule
\textbf{RAG và truy xuất dense}\newline
{\color{schemaGray}Lewis và cs. (2020); Karpukhin và cs. (2020); Asai và cs. (2024); Sarthi và cs. (2024)}
& Ghép bằng chứng truy xuất vào LLM; truy xuất thích ứng, phân cấp.
& Bằng chứng là các đoạn văn độc lập, thiếu cấu trúc suy luận.\\
\addlinespace[1pt]
\textbf{Hỏi đáp đa bước}\newline
{\color{schemaGray}Yang và cs. (2018); Trivedi và cs. (2022)}
& Câu trả lời cần nhiều bằng chứng liên kết.
& Truy xuất riêng lẻ từng đoạn là chưa đủ.\\
\addlinespace[1pt]
\textbf{RAG trên đồ thị và quan hệ nhân quả}\newline
{\color{schemaGray}Edge và cs. (2024); Sun và cs. (2024); Wang và cs. (2025); Parekh và cs. (2025)}
& Tổ chức bằng chứng theo quan hệ; duyệt đồ thị; chuỗi nhân quả.
& Tính đạt được trên đồ thị không tự xác định ngữ nghĩa quy tắc pháp luật.\\
\addlinespace[1pt]
\textbf{Suy luận phản thực và SCM}\newline
{\color{schemaGray}Pearl (2009); Guo và cs. (2023)}
& Mô hình có thể đúng nhờ suy luận rời rạc; can thiệp khác xóa nút.
& Chưa được tích hợp vào pipeline RAG pháp luật.\\
\addlinespace[1pt]
\textbf{Logic ba giá trị, LP đại số tuyến tính}\newline
{\color{schemaGray}Kleene (1952); Fitting (1985); Sato và cs. (2020); Sakama và cs. (2021)}
& Ngữ nghĩa điểm bất động; vector, ma trận, tensor.
& Ít được kết hợp với truy xuất và hỏi đáp pháp luật.\\
\bottomrule
\end{tabular}}

\vspace{0.05cm}
{\scriptsize\color{schemaGray} LP: lập trình logic. Benchmark pháp luật: LexGLUE, CaseHOLD, LegalBench. Danh mục: Phụ lục A7.}
\end{frame}

% ------------------------------------------------------------
% F1.5 Khoảng trống nghiên cứu
% ------------------------------------------------------------
\begin{frame}
\frametitle{Khoảng trống nghiên cứu và vị trí của luận văn}
\begin{columns}[T,onlytextwidth]
\column{0.55\textwidth}
\small
Chưa có framework hỏi đáp pháp luật \textbf{đồng thời}:
\begin{itemize}
  \item truy xuất và xếp hạng đường suy luận nhiều bước;
  \item biểu diễn trạng thái pháp lý chưa đầy đủ bằng ngữ nghĩa ba giá trị;
  \item xử lý phép hội trong một rule và tuyển thay thế giữa nhiều rule;
  \item phân biệt xóa nút tô-pô với can thiệp cấu trúc;
  \item commit có chọn lọc quyết định đã xác minh trước bước sinh.
\end{itemize}

\column{0.41\textwidth}
\centering
\begin{tikzpicture}[font=\scriptsize]
  \def\r{1.05}
  \fill[stepMemory,opacity=0.16] (90:0.68) circle (\r);
  \fill[stepGraph,opacity=0.16] (210:0.68) circle (\r);
  \fill[stepVerify,opacity=0.16] (330:0.68) circle (\r);
  \draw[stepMemory,line width=0.8pt] (90:0.68) circle (\r);
  \draw[stepGraph,line width=0.8pt] (210:0.68) circle (\r);
  \draw[stepVerify,line width=0.8pt] (330:0.68) circle (\r);
  \node[align=center,font=\scriptsize\bfseries,text=stepMemory!75!black] at (0,2.02) {Truy xuất ngữ nghĩa (RAG)};
  \node[align=center,font=\scriptsize,text=stepGraph!75!black] at (-1.5,-1.92) {\textbf{Kết nối cấu trúc}\\(Graph/Causal RAG)};
  \node[align=center,font=\scriptsize,text=stepVerify!85!black] at (1.5,-1.92) {\textbf{Quy tắc hình thức}\\(ba giá trị, SCM)};
  \node[font=\scriptsize\bfseries,text=schemaNavy,fill=white,fill opacity=0.85,text opacity=1,
        rounded corners=2pt,inner sep=1.5pt] at (0,0) {LegalSCM-RAG};
\end{tikzpicture}
\end{columns}

\vspace{0.3cm}
\keymsg{LegalSCM-RAG tích hợp truy xuất đường dẫn nhân quả, suy luận điểm bất động ba giá trị,
xác minh phụ thuộc câu hỏi và cổng commit có kiểm toán.}
\end{frame}

% ------------------------------------------------------------
% F1.6 Câu hỏi nghiên cứu
% ------------------------------------------------------------
\begin{frame}
\frametitle{Bốn câu hỏi nghiên cứu}
\begin{columns}[T,onlytextwidth]
\column{0.49\textwidth}
\rqcard{stepMemory}{RQ1}{Truy xuất có cấu trúc}{Truy xuất trên rule, sự kiện và đường dẫn có tìm
đúng căn cứ pháp lý tốt hơn truy xuất điều luật thô (Vanilla RAG) không?}

\vspace{0.2cm}
\rqcard{stepVerify}{RQ3}{Xác minh chọn lọc}{Xác minh LegalSCM (factual, cạnh trực tiếp,
$do(X_m=0)$) có cải thiện quyết định so với chỉ truy xuất đường dẫn không?
Lợi ích đến từ route nào?}
\column{0.49\textwidth}
\rqcard{stepRetrieve}{RQ2}{Đường suy luận đa bước}{Truy xuất theo đường dẫn nhân quả có chọn
đúng chuỗi suy luận nhiều bước không?}

\vspace{0.2cm}
\rqcard{stepAnswer}{RQ4}{Hiệu quả đầu--cuối}{Câu trả lời và trích dẫn có tốt hơn LLM closed-book,
Vanilla RAG và Causal Path RAG không?}
\end{columns}

\vspace{0.45cm}
\begin{center}
\small\textit{Mỗi câu hỏi được trả lời bằng số liệu và kiểm định ở Phần 5.}
\end{center}
\end{frame}

% ------------------------------------------------------------
% F1.7 Đóng góp
% ------------------------------------------------------------
\begin{frame}
\frametitle{Đóng góp của luận văn}
\begin{columns}[T,onlytextwidth]
\column{0.315\textwidth}
\begin{block}{Dữ liệu và tri thức}
\footnotesize
\begin{itemize}
  \item 2.884 quy tắc điều kiện--hệ quả từ 395 điều BLHS.
  \item Đồ thị nhân quả pháp lý 5.134 nút, 14.420 cạnh.
  \item Benchmark 250 câu hỏi đa bước có gold rule, sự kiện, đường và trích dẫn.
\end{itemize}
\end{block}
\column{0.315\textwidth}
\begin{block}{Phương pháp}
\footnotesize
\begin{itemize}
  \item Truy xuất và xếp hạng đường suy luận đa bước.
  \item Quy tắc ba giá trị dạng đại số tuyến tính: hội, tuyển thay thế, can thiệp cấu trúc.
  \item Xác minh phụ thuộc câu hỏi và cổng commit chọn lọc.
\end{itemize}
\end{block}
\column{0.315\textwidth}
\begin{block}{Đánh giá}
\footnotesize
\begin{itemize}
  \item So sánh có kiểm soát năm cấu hình.
  \item Kiểm định theo cặp: bootstrap, McNemar, hiệu chỉnh Holm.
  \item Kiểm toán route, commit, abstention và can thiệp.
\end{itemize}
\end{block}
\end{columns}

\vspace{0.35cm}
\begin{center}
\kpi[stepVerify]{96,80\%}{Decision Accuracy\\của Full LegalSCM}\hfill
\kpi{+8,0}{điểm so với\\Causal Path RAG}\hfill
\kpi{0,0011}{giá trị $p$\\(hiệu chỉnh Holm)}\hfill
\kpi{250}{câu hỏi BLHS\\đa bước}
\end{center}
\end{frame}
